\IfFileExists{stacks-project.cls}{%
\documentclass{stacks-project}
}{%
\documentclass{amsart}
}
\usepackage{amssymb}

% For dealing with references we use the comment environment
\usepackage{verbatim}
\usepackage{fancyvrb}
\newenvironment{reference}{\comment}{\endcomment}
%\newenvironment{reference}{}{}
\newenvironment{slogan}{\comment}{\endcomment}
\newenvironment{history}{\comment}{\endcomment}

% For source figures retained as assets
\usepackage{graphicx}

% For commutative diagrams we use Xy-pic
\usepackage[all]{xy}

% We use 2cell for 2-commutative diagrams.
\xyoption{2cell}
\UseAllTwocells

% We use multicol for the list of chapters between chapters
\usepackage{multicol}

% This is generally recommended for better output
\usepackage{lmodern}
\usepackage[T1]{fontenc}
\usepackage{textalpha}

% For cross-file-references
% We prefer the snapshot tmp/xref/#2.aux,
% which latexmkrc replaces atomically, so
% that documents compiled simultaneously
% never read a half-written .aux file.
% The last argument keeps links to #2.pdf.
\usepackage{xr-hyper}
\newcommand{\maybeexternaldocument}[2]{%
  \IfFileExists{#2.aux}{%
    \externaldocument[#1]{#2}%
  }{%
    \PackageWarning{stack}{Missing #2.aux; external references with prefix #1 unavailable}%
  }%
}
\let\externaldocumentorig\externaldocument
\renewcommand{\externaldocument}[2][]{%
  \IfFileExists{tmp/xref/#2.aux}{%
    \externaldocumentorig[#1]{tmp/xref/#2}[#2.pdf]%
  }{%
  \IfFileExists{#2.aux}{%
    \externaldocumentorig[#1]{#2}%
  }{%
    \PackageWarning{stack}{Missing #2.aux; external references with prefix #1 unavailable}%
  }%
  }%
}

% Package for hypertext links:
\usepackage{hyperref}

% For any local file, say "hello.tex" you want to link to please
% use \externaldocument[hello-]{hello}
\externaldocument[introduction-]{introduction}
\externaldocument[conventions-]{conventions}
\externaldocument[sets-]{sets}
\externaldocument[categories-]{categories}
\externaldocument[topology-]{topology}
\externaldocument[sheaves-]{sheaves}
\externaldocument[sites-]{sites}
\externaldocument[stacks-]{stacks}
\externaldocument[fields-]{fields}
\externaldocument[algebra-]{algebra}
\externaldocument[brauer-]{brauer}
\externaldocument[homology-]{homology}
\externaldocument[derived-]{derived}
\externaldocument[simplicial-]{simplicial}
\externaldocument[more-algebra-]{more-algebra}
\externaldocument[smoothing-]{smoothing}
\externaldocument[modules-]{modules}
\externaldocument[sites-modules-]{sites-modules}
\externaldocument[injectives-]{injectives}
\externaldocument[cohomology-]{cohomology}
\externaldocument[sites-cohomology-]{sites-cohomology}
\externaldocument[dga-]{dga}
\externaldocument[dpa-]{dpa}
\externaldocument[sdga-]{sdga}
\externaldocument[hypercovering-]{hypercovering}
\externaldocument[schemes-]{schemes}
\externaldocument[constructions-]{constructions}
\externaldocument[properties-]{properties}
\externaldocument[morphisms-]{morphisms}
\externaldocument[coherent-]{coherent}
\externaldocument[divisors-]{divisors}
\externaldocument[limits-]{limits}
\externaldocument[varieties-]{varieties}
\externaldocument[topologies-]{topologies}
\externaldocument[descent-]{descent}
\externaldocument[perfect-]{perfect}
\externaldocument[more-morphisms-]{more-morphisms}
\externaldocument[flat-]{flat}
\externaldocument[groupoids-]{groupoids}
\externaldocument[more-groupoids-]{more-groupoids}
\externaldocument[etale-]{etale}
\externaldocument[chow-]{chow}
\externaldocument[intersection-]{intersection}
\externaldocument[pic-]{pic}
\externaldocument[weil-]{weil}
\externaldocument[adequate-]{adequate}
\externaldocument[dualizing-]{dualizing}
\externaldocument[duality-]{duality}
\externaldocument[discriminant-]{discriminant}
\externaldocument[derham-]{derham}
\externaldocument[local-cohomology-]{local-cohomology}
\externaldocument[algebraization-]{algebraization}
\externaldocument[curves-]{curves}
\externaldocument[resolve-]{resolve}
\externaldocument[models-]{models}
\externaldocument[functors-]{functors}
\externaldocument[equiv-]{equiv}
\externaldocument[pione-]{pione}
\externaldocument[etale-cohomology-]{etale-cohomology}
\externaldocument[proetale-]{proetale}
\externaldocument[relative-cycles-]{relative-cycles}
\externaldocument[more-etale-]{more-etale}
\externaldocument[trace-]{trace}
\externaldocument[crystalline-]{crystalline}
\externaldocument[spaces-]{spaces}
\externaldocument[spaces-properties-]{spaces-properties}
\externaldocument[spaces-morphisms-]{spaces-morphisms}
\externaldocument[decent-spaces-]{decent-spaces}
\externaldocument[spaces-cohomology-]{spaces-cohomology}
\externaldocument[spaces-limits-]{spaces-limits}
\externaldocument[spaces-divisors-]{spaces-divisors}
\externaldocument[spaces-over-fields-]{spaces-over-fields}
\externaldocument[spaces-topologies-]{spaces-topologies}
\externaldocument[spaces-descent-]{spaces-descent}
\externaldocument[spaces-perfect-]{spaces-perfect}
\externaldocument[spaces-more-morphisms-]{spaces-more-morphisms}
\externaldocument[spaces-flat-]{spaces-flat}
\externaldocument[spaces-groupoids-]{spaces-groupoids}
\externaldocument[spaces-more-groupoids-]{spaces-more-groupoids}
\externaldocument[bootstrap-]{bootstrap}
\externaldocument[spaces-pushouts-]{spaces-pushouts}
\externaldocument[spaces-chow-]{spaces-chow}
\externaldocument[groupoids-quotients-]{groupoids-quotients}
\externaldocument[spaces-more-cohomology-]{spaces-more-cohomology}
\externaldocument[spaces-simplicial-]{spaces-simplicial}
\externaldocument[spaces-duality-]{spaces-duality}
\externaldocument[formal-spaces-]{formal-spaces}
\externaldocument[restricted-]{restricted}
\externaldocument[spaces-resolve-]{spaces-resolve}
\externaldocument[formal-defos-]{formal-defos}
\externaldocument[defos-]{defos}
\externaldocument[cotangent-]{cotangent}
\externaldocument[examples-defos-]{examples-defos}
\externaldocument[algebraic-]{algebraic}
\externaldocument[examples-stacks-]{examples-stacks}
\externaldocument[stacks-sheaves-]{stacks-sheaves}
\externaldocument[criteria-]{criteria}
\externaldocument[artin-]{artin}
\externaldocument[quot-]{quot}
\externaldocument[stacks-properties-]{stacks-properties}
\externaldocument[stacks-morphisms-]{stacks-morphisms}
\externaldocument[stacks-limits-]{stacks-limits}
\externaldocument[stacks-cohomology-]{stacks-cohomology}
\externaldocument[stacks-perfect-]{stacks-perfect}
\externaldocument[stacks-introduction-]{stacks-introduction}
\externaldocument[stacks-more-morphisms-]{stacks-more-morphisms}
\externaldocument[stacks-geometry-]{stacks-geometry}
\externaldocument[moduli-]{moduli}
\externaldocument[moduli-curves-]{moduli-curves}
\externaldocument[examples-]{examples}
\externaldocument[exercises-]{exercises}
\externaldocument[guide-]{guide}
\externaldocument[desirables-]{desirables}
\externaldocument[coding-]{coding}
\externaldocument[obsolete-]{obsolete}
\externaldocument[fdl-]{fdl}
\externaldocument[index-]{index}

% Heap Project documents
\externaldocument[k3-]{k3}
\externaldocument[torelli-]{torelli}
\externaldocument[linear-algebra-]{linear-algebra}
\externaldocument[groups-]{groups}
\externaldocument[complex-analysis-]{complex-analysis}
\externaldocument[functional-analysis-]{functional-analysis}
\externaldocument[categories-]{categories}
\externaldocument[commutative-algebra-]{commutative-algebra}
\externaldocument[ringed-spaces-]{ringed-spaces}
\externaldocument[homological-algebra-]{homological-algebra}
\externaldocument[lie-algebras-]{lie-algebras}
\externaldocument[representation-theory-]{representation-theory}
\externaldocument[smooth-manifolds-]{smooth-manifolds}
\externaldocument[differential-topology-]{differential-topology}
\externaldocument[algebraic-topology-]{algebraic-topology}
\externaldocument[differential-geometry-]{differential-geometry}
\externaldocument[algebraic-geometry-]{algebraic-geometry}
\externaldocument[symplectic-geometry-]{symplectic-geometry}
\externaldocument[complex-geometry-]{complex-geometry}
\externaldocument[hodge-theory-]{hodge-theory}
\externaldocument[%
variations-of-hodge-structures-]{%
variations-of-hodge-structures}
\externaldocument[%
homological-mirror-symmetry-]{%
homological-mirror-symmetry}
\maybeexternaldocument{deformations-}{deformations}
\externaldocument[vertex-algebras-]{vertex-algebras}
\externaldocument[springer-]{springer}
\externaldocument[infty-categories-]{infty-categories}
\externaldocument[seminars-differential-geometry-]{%
seminars-differential-geometry}
\externaldocument[seminars-complex-geometry-]{%
seminars-complex-geometry}
\externaldocument[seminars-algebraic-geometry-]{%
seminars-algebraic-geometry}
\externaldocument[seminars-others-]{%
seminars-others}
\externaldocument[%
geometric-prequantization-]{%
geometric-prequantization}
\externaldocument[surfaces-]{surfaces}
\externaldocument[birational-maps-commutative-algebra-]{birational-maps-commutative-algebra}
\externaldocument[cremona-transformations-]{cremona-transformations}
\externaldocument[derived-categories-of-sheaves-]{derived-categories-of-sheaves}
\externaldocument[geometric-invariant-theory-]{geometric-invariant-theory}
\externaldocument[geometric-stability-conditions-and-group-actions-]{geometric-stability-conditions-and-group-actions}
\externaldocument[moduli-spaces-of-sheaves-]{moduli-spaces-of-sheaves}
\externaldocument[bridgeland-stability-general-theory-]{bridgeland-stability-general-theory}
\externaldocument[hilbert-schemes-]{hilbert-schemes}
\externaldocument[noncommutative-abelian-surfaces-and-kummer-type-hyperkahler-varieties-]{noncommutative-abelian-surfaces-and-kummer-type-hyperkahler-varieties}
\externaldocument[mumford-tate-groups-in-hodge-theory-]{mumford-tate-groups-in-hodge-theory}
\externaldocument[hodge-birational-atoms-2026-]{hodge-birational-atoms-2026}
\externaldocument[xxii-egd-2026-]{%
xxii-egd-2026}
\externaldocument[pde-]{pde}
\externaldocument[probability-]{probability}

\externaldocument[linguistics-]{linguistics}
\externaldocument[genealogia-familiar-]{%
genealogia-familiar}

\externaldocument[%
16th-alga-meeting-2026-]{%
16th-alga-meeting-2026}

\externaldocument[%
iv-jornadas-lefschetz-2026-]{%
iv-jornadas-lefschetz-2026}

\externaldocument[reading-the-masters-]{%
reading-the-masters}

\externaldocument[real-analysis-]{%
real-analysis}

\externaldocument[optimal-transport-]{%
optimal-transport}

\externaldocument[history-]{history}

\externaldocument[seiberg-witten-]{seiberg-witten}

\externaldocument[history-philosophy-biology-]{history-philosophy-biology}

\externaldocument[chemistry-]{chemistry}

\externaldocument[molecular-cell-biology-]{molecular-cell-biology}

\externaldocument[microbiology-]{microbiology}

\externaldocument[protistology-]{protistology}

\externaldocument[mycology-]{mycology}

\externaldocument[botany-]{botany}

\externaldocument[zoology-]{zoology}

\externaldocument[animal-histology-]{animal-histology}

\externaldocument[animal-physiology-]{animal-physiology}

\externaldocument[systematics-biogeography-]{systematics-biogeography}

\externaldocument[ecology-natural-resources-]{ecology-natural-resources}

\externaldocument[earth-sciences-]{earth-sciences}

\externaldocument[biostatistics-]{biostatistics}

\externaldocument[mathematical-systems-biology-]{mathematical-systems-biology}

\externaldocument[cimpa-el-salvador-2024-]{cimpa-el-salvador-2024}

% Theorem environments.
%
\theoremstyle{plain}
\newtheorem{theorem}[subsection]{Theorem}
\newtheorem{proposition}[subsection]{Proposition}
\newtheorem{lemma}[subsection]{Lemma}
\newtheorem{corollary}[subsection]{Corollary}

\theoremstyle{definition}
\newtheorem{definition}[subsection]{Definition}
\newtheorem{example}[subsection]{Example}
\newtheorem{exercise}[subsection]{Exercise}
\newtheorem{situation}[subsection]{Situation}

\theoremstyle{remark}
\newtheorem{remark}[subsection]{Remark}
\newtheorem{remarks}[subsection]{Remarks}

\numberwithin{equation}{subsection}

% Macros
%
\def\lim{\mathop{\mathrm{lim}}\nolimits}
\def\colim{\mathop{\mathrm{colim}}\nolimits}
\def\Spec{\mathop{\mathrm{Spec}}}
\def\Hom{\mathop{\mathrm{Hom}}\nolimits}
\def\Ext{\mathop{\mathrm{Ext}}\nolimits}
\def\SheafHom{\mathop{\mathcal{H}\!\mathit{om}}\nolimits}
\def\SheafExt{\mathop{\mathcal{E}\!\mathit{xt}}\nolimits}
\def\Sch{\mathit{Sch}}
\def\Mor{\mathop{\mathrm{Mor}}\nolimits}
\def\Ob{\mathop{\mathrm{Ob}}\nolimits}
\def\Sh{\mathop{\mathit{Sh}}\nolimits}
\def\NL{\mathop{N\!L}\nolimits}
\def\CH{\mathop{\mathrm{CH}}\nolimits}
\def\proetale{{pro\text{-}\acute{e}tale}}
\def\etale{{\acute{e}tale}}
\def\QCoh{\mathit{QCoh}}
\def\Ker{\mathop{\mathrm{Ker}}}
\def\Im{\mathop{\mathrm{Im}}}
\def\Coker{\mathop{\mathrm{Coker}}}
\def\Coim{\mathop{\mathrm{Coim}}}

% Boxtimes
%
\DeclareMathSymbol{\boxtimes}{\mathbin}{AMSa}{"02}

%
% Macros for moduli stacks/spaces
%
\def\QCohstack{\mathcal{QC}\!\mathit{oh}}
\def\Cohstack{\mathcal{C}\!\mathit{oh}}
\def\Spacesstack{\mathcal{S}\!\mathit{paces}}
\def\Quotfunctor{\mathrm{Quot}}
\def\Hilbfunctor{\mathrm{Hilb}}
\def\Curvesstack{\mathcal{C}\!\mathit{urves}}
\def\Polarizedstack{\mathcal{P}\!\mathit{olarized}}
\def\Complexesstack{\mathcal{C}\!\mathit{omplexes}}
% \Pic is the operator that assigns to X its picard group, usage \Pic(X)
% \Picardstack_{X/B} denotes the Picard stack of X over B
% \Picardfunctor_{X/B} denotes the Picard functor of X over B
\def\Pic{\mathop{\mathrm{Pic}}\nolimits}
\def\Picardstack{\mathcal{P}\!\mathit{ic}}
\def\Picardfunctor{\mathrm{Pic}}
\def\Deformationcategory{\mathcal{D}\!\mathit{ef}}


\begin{document}

\title{Ecology and Natural Resources}
\maketitle

\phantomsection
\label{section-phantom}

\tableofcontents

\section{What this chapter covers}
\label{section-overview}

A short chapter: the basic vocabulary of
ecology (populations, communities,
ecosystems, biomes, energy flow and nutrient
cycles), the equilibrium theory of island
biogeography, and then conservation and
natural resources as taught in a Mexican
course: endangered-species lists, Mexican
environmental law, energy resources and
peak oil, and two critical readings on
society and nature.

\textbf{Sources of these notes.}
Daniel took the one-semester course
\emph{Recursos Naturales} (Facultad de
Ciencias, UNAM, semester 2015-2). Its only
surviving trace is his answered third
partial exam, on the IUCN Red List, CITES,
the Mexican norm NOM-059, environmental
impact statements, ecological land-use
planning, obsolescence, and Hubbert's peak
oil curve with EROI. From the
\emph{Biogeograf\'ia} course (third
semester, autumn 2014) come his exam notes
on island biogeography and his review of a
climate-change documentary. The basic
ecology is written from Stern et al.'s
\emph{Introductory Plant Biology}
\cite{stern2008plant} and Hickman et
al., \emph{Animal Diversity}
\cite{hickman2002animal}; island
biogeography is checked against Hubbell
\cite{hubbell2001unified}; conservation
genetics against Freeland
\cite{freeland2005molecular}; and the
society--nature section rests on Galafassi
\cite{galafassi2000sociedad} and Levins
and Lewontin \cite{levins1985dialectical}.

\section{Populations, communities and biomes}
\label{section-levels}

This section is written from the textbooks;
the surviving coursework has no material on
it.

\begin{definition}
\label{definition-levels}
A \emph{population} is a group of
individuals of one species in an area. A
\emph{community} (biotic community) is the
set of all populations of living organisms
in an area. An \emph{ecosystem} is a
community together with its physical
environment, linked by physical, chemical
and biological processes. A \emph{biome} is
a terrestrial community considered on a
continental or global scale.
\end{definition}

These definitions follow Stern
\cite[Chapters 25 and 26]{stern2008plant}.
A population is described by its numbers,
density (individuals per unit area or
volume), genetic diversity and biomass. The
distribution of a species depends on abiotic
factors (light, temperature, precipitation,
soil, fire) and on biotic factors
(competition, grazing). Climate is the main
factor deciding the composition of a biome;
the biomes of North America are tundra,
taiga, temperate deciduous forest,
grassland, desert, mountain and coastal
forest, and tropical rain forest
\cite[Chapter 26]{stern2008plant}.

\textbf{Population growth.}
Write $N$ for population size and $r$ for
the intrinsic rate of increase per
individual (birth rate minus death rate).
With unlimited resources growth is
exponential; when a limiting resource caps
the population at the carrying capacity
$K$, growth follows the logistic equation
\cite[Chapter 2]{hickman2002animal}:
$$
\frac{dN}{dt} = rN, \qquad
\frac{dN}{dt} = rN\,\frac{K-N}{K}.
$$
The logistic curve is sigmoid: growth slows
as $N$ approaches $K$ and stops at $N=K$.
This negative feedback of density on growth
is density dependence. Real populations
fluctuate around $K$, because $K$ changes and
because populations respond with a lag.

\textbf{Interactions.}
Ecologists mark the effect of an
interaction on each partner as beneficial
($+$), harmful ($-$) or neutral ($0$):
predation, parasitism and herbivory are
$+-$, commensalism is $+0$, mutualism is
$++$, and competition harms both. The
\emph{niche} of a species is its way of life
in the community; by the competitive
exclusion principle no two species occupy
the same niche indefinitely. A
\emph{keystone species} is one whose removal
changes the whole community, like the sea
star \emph{Pisaster ochraceus} on Pacific
rocky shores
\cite[Chapter 2]{hickman2002animal}.

\section{Energy flow and nutrient cycles}
\label{section-energy-cycles}

This section is written from the textbooks;
the surviving coursework has no material on
it.

\textbf{Energy flow.}
Producers (plants, algae) capture light by
photosynthesis; primary consumers eat
producers, secondary consumers eat primary
consumers; decomposers, mainly bacteria and
fungi, return organic matter to forms that
producers can use. Feeding links form food
chains and food webs, rarely with more than
six links. Energy flows one way and cannot
be recycled: only about 1\% of the light
reaching a temperate community enters
organic matter, and only about 10\% of the
energy in one link passes to the next, the
rest being lost as heat or left unused.
Hence the pyramid of energy, and the
greater efficiency of eating low in the
food chain
\cite[Chapter 25]{stern2008plant}.

\textbf{Nutrient cycles.}
Matter, unlike energy, cycles between
organic and inorganic forms
\cite[Chapter 25]{stern2008plant}:
\begin{itemize}
\item Water evaporates from water bodies and
is transpired by plants, condenses, and
returns as precipitation; some sinks to the
water table and the aquifers.
\item Carbon is fixed from $\mathrm{CO_2}$ by
photosynthesis and returned by respiration,
most of it by decomposing bacteria and
fungi; burning fossil fuels and forests
adds more; anaerobic bacteria release
methane.
\item Nitrogen makes up about 78\% of the
atmosphere but $\mathrm{N_2}$ is inert.
Nitrogen-fixing bacteria, free in the soil
or in the root nodules of legumes, convert
it to ammonia; decomposers release nitrogen
from dead tissue; denitrifying bacteria in
poorly aerated soils use nitrate in place of
oxygen and return nitrogen to the air;
lightning fixes a little. Harvests, erosion
and fire remove nitrogen from soils.
\end{itemize}

\section{Island biogeography}
\label{section-island-biogeography}

The biogeography course presented the
equilibrium theory of island biogeography
of MacArthur and Wilson. An ``island'' is any
patch of habitat surrounded by very
different habitat. The number of species on
an island is a dynamic equilibrium $S^*$
between the rate $I$ of immigration of new
species from a source area and the rate $E$
of extinction of resident species: as the
number $S$ of species on the island grows,
$I$ falls (fewer arrivals are new) and $E$
rises, and the curves cross at $S^*$
\cite[Chapter 1]{hubbell2001unified}. Larger
islands have lower extinction rates and so
more species; more distant islands have
lower immigration rates and so fewer
species. The equilibrium fixes the number,
not the identity, of species: all species
are treated as identical, which makes the
theory a neutral theory and the starting
point of Hubbell's own
\cite[Chapter 1]{hubbell2001unified}.

Daniel's exam notes list the criticisms
discussed in class: species can arise on the
island itself, islands are not homogeneous,
colonization and extinction are not
independent, and history is ignored. Island
theory was later applied to the design of
nature reserves
\cite[Chapter 1]{hubbell2001unified}.

\section{Species at risk}
\label{section-species-at-risk}

\textbf{The IUCN Red List.}
The International Union for Conservation of
Nature (IUCN) compiles Red Lists of species
at risk; the categories rest on parameters
such as current population size
\cite[Chapter 7]{freeland2005molecular}. The
nine categories, as listed in Daniel's exam,
are: Extinct, Extinct in the Wild,
Critically Endangered, Endangered,
Vulnerable, Near Threatened, Least Concern,
Data Deficient and Not Evaluated. The three
categories Critically Endangered, Endangered
and Vulnerable together form the
``threatened'' species. Small populations are
at risk not only from chance events but also
because genetic drift erodes their genetic
diversity and inbreeding becomes common
\cite[Chapter 7]{freeland2005molecular}.

\textbf{CITES.}
The Convention on International Trade in
Endangered Species of Wild Fauna and Flora
(in force since 1975) lists species in three
appendices:
\begin{itemize}
\item Appendix I: species threatened with
extinction; commercial trade in wild
specimens is forbidden;
\item Appendix II: species not necessarily
threatened now, whose trade must be
controlled so that it does not endanger
them (the longest list);
\item Appendix III: species protected in at
least one country that has asked the other
parties for help in controlling their
trade.
\end{itemize}

\textbf{NOM-059.}
The Mexican official norm NOM-059 lists the
species at risk in Mexico. As in the course
(2001 version), its categories are: in
danger of extinction (populations or range
greatly reduced), threatened (would become
endangered if the causes of decline
continue), and subject to special
protection (would become threatened without
protection). A species is placed by the
Method for Evaluating the Risk of
Extinction (MER), which scores four
criteria:
\begin{enumerate}
\item the extent of its distribution in
Mexico;
\item the state of its habitat;
\item its intrinsic biological
vulnerability;
\item the impact of human activity on it.
\end{enumerate}
The criteria are scored numerically, with
different scales for animals and plants.

\section{Environmental planning in Mexico}
\label{section-planning}

These tools are described as in the course.

\textbf{Environmental impact statement.}
A \emph{Manifestaci\'on de Impacto Ambiental}
(MIA) is a study, made before a project is
authorized, that estimates its effects on
the environment and proposes alternatives
or mitigation measures, so that the
authority can decide whether to allow it.
It is required for public works, hydraulic
works, industry, mineral processing,
tourism developments and housing
developments, among others.

\textbf{Ecological land-use planning.}
The \emph{ordenamiento ecol\'ogico} seeks a
balance between the natural aptitude of the
land, human settlements, productive
activities and natural hazards. The exam
also asked about multicriteria analysis of
conflicts between land users; Daniel's
answer is blank.

\textbf{Obsolescence.}
Perceived obsolescence is the urge to
replace a product that still works;
programmed (planned) obsolescence is the
design of products to fail early.

\section{Energy resources}
\label{section-energy}

This section is described as in the course.

\textbf{Hubbert's curve.}
Hubbert's gross curve models the yearly
production of oil from a region as a
symmetric bell: production rises to a
maximum, the peak, and falls afterwards,
with about as much oil produced after the
peak as before it.

\textbf{EROI.}
The energy return on investment is the
ratio of the energy obtained to the energy
spent obtaining it. Writing $P$ for the
energy produced, $C$ for the energy
invested and $R$ for the EROI,
$$
R = \frac{P}{C}, \qquad
P - C = P\left(1 - \frac{1}{R}\right).
$$
The net energy available to society is the
second quantity. Since the easiest oil is
extracted first, $R$ falls with time, and
the net Hubbert curve lies further and
further below the gross one after the peak:
less usable oil remains than the gross curve
suggests.

\textbf{Alternatives.}
Wind, solar photovoltaic and solar thermal,
hydroelectric, tidal, geothermal energy and
biofuels are renewable, but Daniel's answer
argues that none can replace oil in the
time oil has left: each needs large and
costly installations, depends on
uncontrollable natural conditions, or has
its own environmental impact, and none
matches the amount of energy oil provides.

\textbf{Climate change.}
In the biogeography course Daniel reviewed
the documentary \emph{Six Degrees Could
Change the World}. As in the documentary:
the present climate is almost one degree
warmer than in previous centuries, the
warming is caused by human emissions of
$\mathrm{CO_2}$, a two-degree rise is a
critical point because positive feedbacks
(warming causes drought, which causes more
warming) may then put it beyond control, and
past climate changes were much slower than
the present one.

\section{Society and nature}
\label{section-society-nature}

Two readings question how ecology treats
human society.

Galafassi argues that environmental
problems are social in origin: they arise
from the way a society appropriates nature
to build its habitat and to produce. The
natural and the social are different levels
of reality, each with its own laws, so
neither can be reduced to the other; the
ecosystem view, which adds humans as one
more component, misses this. He proposes to
study their articulation through the
production process, and in particular
through labour and technology: technology
is not an independent force but is chosen
by economic rationality, and choices that
ignore the rhythms of the environment
degrade it. Agriculture is his main
example, since there labour acts directly on
living matter \cite{galafassi2000sociedad}.

Levins and Lewontin attack the dichotomy
between mechanistic reductionism and
holism in ecology. Reductionism treats each
species as reacting to an external
environment and forgets that organisms also
change their environment (plants change the
soil they grow in). Holism, as in Clements's
community as a superorganism, subordinates
the parts to an abstract balance of nature.
Their example: in a predator--prey system
each species alone oscillates
unpredictably, while the pair moves
regularly along a spiral towards
equilibrium, a property of the coupled
system and not of either species
\cite[Chapter 6]{levins1985dialectical}.

\bibliography{my}
\bibliographystyle{amsalpha}

\end{document}
